\documentclass[11pt]{article}

\usepackage[a4paper,margin=25mm]{geometry}
\usepackage{amsmath,amssymb,mathtools}
\usepackage{booktabs,longtable,array}
\usepackage[hidelinks]{hyperref}
\usepackage[T1]{fontenc}

\newcommand{\R}{\mathbb{R}}
\newcommand{\C}{\mathcal{C}}
\newcommand{\Sset}{\mathcal{S}}
\newcommand{\norm}[1]{\left\lVert #1\right\rVert}
\newcommand{\abs}[1]{\left\lvert #1\right\rvert}
\newcommand{\kernel}[1]{\nolinkurl{#1}}
\newcommand{\code}[1]{\nolinkurl{#1}}

\title{Active-Row Subclade Pruning in the \texttt{gpurec} Adjoint}
\author{gpurec developers}
\date{\today}

\begin{document}
\maketitle
\tableofcontents

\section{Scope and conclusion}

The flags \code{USE_ACTIVE} and \code{USE_ACTIVE_MASK} specialize kernels to
skip gene-clade rows whose incoming adjoint is zero or sufficiently small.
Despite their short names, they do \emph{not} mask species-tree states.  There
is one Boolean per gene-clade row in the current wave, and that Boolean applies
to all species states of the row.

The mathematical reason the optimization works is row separability.  After
the forward reconciliation has been fixed, every clade row has its own
within-row adjoint system
\begin{equation}
  (I-J_i^T)v_i=b_i.
  \label{eq:local-adjoint}
\end{equation}
If the accumulated right-hand side $b_i$ is exactly zero and the system has a
unique solution, then $v_i=0$.  Consequently the row contributes neither
parameter cotangents nor cotangents to its child clades.  It is therefore exact
to omit that row and to stop reverse propagation through that route.  Applied
in reverse topological order, this can remove a whole descendant subclade.

There is an important qualification.  A positive pruning threshold also drops
\emph{small nonzero} $b_i$.  That is a deliberate numerical approximation, not
an algebraic identity.  Exact-zero pruning is exact; threshold pruning is
correct as an implementation of the explicitly truncated adjoint and is close
to the full gradient only when the discarded cotangents are not strongly
amplified by Equation~\eqref{eq:local-adjoint}.  Section~\ref{sec:error}
states the corresponding bound and assumptions.

\section{Clade waves and index domains}

Let $\C$ be the gene-clade directed acyclic graph.  A recorded split of parent
clade $i$ has child clades $\ell$ and $r$, so the forward dependency edges are
\begin{equation}
  (i,\ell),\ (i,r):\qquad
  \Pi_i \text{ depends on } \Pi_\ell \text{ and } \Pi_r.
\end{equation}
Preprocessing orders clades from leaves toward roots and groups equal-level
rows into waves.  For a wave beginning at global row \code{ws} and containing
$W$ rows, the three relevant domains are
\begin{align}
  w &\in \{0,\ldots,W-1\}, &&\text{wave-local clade row},\\
  i &= \mathtt{ws}+w, &&\text{global clade row},\\
  s &\in \{0,\ldots,S-1\}, &&\text{species-tree state}.
\end{align}
For split $n$, \code{reduce\_idx[n]} is its wave-local parent $w_n$; the child
arrays \code{sl[n]} and \code{sr[n]} contain global clade rows.  Thus the
activity array has shape $[W]$ and is read as
\begin{equation}
  a_{w_n}=\mathtt{active\_mask[reduce\_idx[n]]}.
  \label{eq:index-contract}
\end{equation}
It is never indexed by $s$, and it must not be confused with the separate
split-side mask of shape $[2N_{\rm split}]$ discussed in
Section~\ref{sec:side-mask}.

Rows in one wave have no cross-row fixed-point dependence: their children are
in earlier waves, while their self-loop is entirely along the species axis of
the same row.  This makes the within-wave Jacobian block diagonal,
\begin{equation}
  J_{\rm wave}=\operatorname{diag}(J_0,\ldots,J_{W-1}),
  \label{eq:block-diagonal}
\end{equation}
which is the structural fact that permits a single Boolean to remove a whole
row from every stage of its local solve.

\section{How the active mask is constructed}

At the start of the reverse pass, root rows receive the likelihood cotangent.
All other entries of the dense array \code{accumulated\_rhs} are initialized to
zero.  Processing waves from roots toward leaves atomically scatters child
cotangents into that array.  Immediately before wave $k$ is processed, its
slice
\begin{equation}
  b_{w,s}=\mathtt{accumulated\_rhs[ws+w,s]}
  \label{eq:rhs}
\end{equation}
therefore contains the \emph{sum of all reverse paths already reaching that
clade}.  Activity is based on the row infinity norm
\begin{equation}
  r_w=\norm{b_w}_\infty=\max_s\abs{b_{w,s}}.
  \label{eq:row-max}
\end{equation}
The kernel \kernel{_select_active_adjoint_rows_kernel} scans all $S$ entries
and writes one Boolean.  The current host helper selects one of two boundary
conventions:
\begin{equation}
  a_w =
  \begin{cases}
    [r_w\ge \tau], & \mathtt{use\_pruning=True},\\
    [r_w>\tau], & \mathtt{use\_pruning=False},
  \end{cases}
  \label{eq:actual-mask}
\end{equation}
where $[\cdot]$ is an indicator and $\tau$ is
\code{adjoint\_pruning\_threshold}.  The distinction is only equality at the
threshold.  In particular, the current \code{use\_adjoint\_pruning=False}
path does not make every row active and does not remove the mask; at positive
$\tau$ it still drops rows with $r_w\le\tau$.  To request the exact-zero
criterion with the current interface, use $\tau=0$ and the strict $>$ branch.
To force all rows active through this helper, use $\tau=0$ and the $\ge$
branch.  This detail is an interface fact, not part of the mathematical
justification for pruning.

Because Equation~\eqref{eq:row-max} uses the accumulated \emph{net}
cotangent, exact cancellation between reverse paths may make a row inactive.
That is valid: reverse-mode differentiation is linear in cotangents, so a zero
net cotangent has a zero downstream VJP even when its individual summands were
nonzero.

\section{The reverse recurrence and subclade pruning}

For each row $i$, write the converged within-row recurrence abstractly as
\begin{equation}
  x_i=F_i(x_i,\{x_j:j\in\operatorname{ch}(i)\},p),
  \label{eq:row-map}
\end{equation}
where $x_i\in\R^S$ is the row of $\Pi$ values, $p$ denotes model parameters,
and $\operatorname{ch}(i)$ contains the clades appearing as children in the
DTS split terms.  At the fixed primal point define
\begin{align}
  J_i &= \frac{\partial F_i}{\partial x_i}, &
  C_{ij} &= \frac{\partial F_i}{\partial x_j}, &
  B_i &= \frac{\partial F_i}{\partial p}.
\end{align}
If $b_i$ is the cotangent accumulated from the objective and all already
processed parents, implicit reverse differentiation gives
\begin{align}
  v_i &= (I-J_i^T)^{-1}b_i, \label{eq:solve}\\
  b_j &\mathrel{+}= C_{ij}^T v_i
    && (j\in\operatorname{ch}(i)), \label{eq:child-scatter}\\
  \nabla_p L &\mathrel{+}= B_i^T v_i. \label{eq:param-scatter}
\end{align}
The actual kernels decompose $C_{ij}^Tv_i$ and $B_i^Tv_i$ into duplication,
transfer, speciation, extinction, receiver-weight, and leaf terms, but all are
linear in the solved cotangent $v_i$.

If $b_i=0$, Equation~\eqref{eq:solve} gives $v_i=0$.  Equations
\eqref{eq:child-scatter}--\eqref{eq:param-scatter} then give zero.  Skipping
row $i$ is consequently equivalent to evaluating all three equations.  Now
process the waves in reverse topological order.  A child reached only through
inactive parents receives no cotangent, is itself classified inactive, and
the argument repeats.  This is the precise sense in which the optimization
prunes a subclade.

The reconciliation structure is a clade DAG, not necessarily a tree.  A child
may receive cotangents from several parent splits.  Inactivating one parent
only removes that route.  If another active route supplies a sufficiently
large cotangent, the child row remains active.  Thus ``pruned subclade'' means
the portion of the reverse dependency DAG no longer reachable through active
cotangent edges, not blindly every descendant of one inactive split.

\section{What each masked stage does}

\subsection{Backward-only DTS recomputation: \texttt{USE\_ACTIVE}}

The ordinary primal forward pass calls \code{compute\_dts\_forward} without an
activity array.  During the adjoint, DTS values are recomputed for the current
wave with \code{active\_parent\_rows=active\_mask}.  This makes
\code{USE\_ACTIVE} a compile-time true specialization in three kernels:
\begin{itemize}
  \item \kernel{_reduce_single_gene_split_events_kernel} returns for an inactive parent after writing
        its output row to $-\infty$ and its gauge offset to zero;
  \item \kernel{_stage_multiple_gene_split_event_reduction_kernel} does not form partial log-sum-exp
        tiles for an inactive parent;
  \item \kernel{_finalize_multiple_gene_split_event_reduction_kernel} writes the inactive output row to
        $-\infty$ without reading the uninitialized stage-1 partials.
\end{itemize}
The output tensor was initialized to $-\infty$ and its offset tensor to zero.
More importantly, every later backward load for that row is also guarded by
the same activity bit.  The sentinel values therefore establish a valid
inactive representation; they are not interpreted as a new model value.

\subsection{Within-row adjoint: \texttt{USE\_ACTIVE\_MASK}}

The retained wave backward path solves Equation~\eqref{eq:local-adjoint} by a
Neumann series or GMRES.  The mask is applied consistently to every component:
\begin{enumerate}
  \item \kernel{_prepare_reconciliation_self_loop_vjp_kernel} forms event
        probabilities and the coefficients of $J_i^T$ only for active rows;
  \item \kernel{_apply_reconciliation_self_loop_transpose_kernel} applies $J_i^T$ only on
        active rows during every Neumann or GMRES iteration;
  \item the transfer receiver-log-weight VJP kernel omits inactive
        receiver-log-weight cotangents;
  \item \kernel{_accumulate_reconciliation_event_vjp_kernel} contracts $v_i$
        with event probabilities only for active rows and writes zero to its
        per-row gradient outputs for inactive rows.
\end{enumerate}
For GMRES, the host also explicitly forms
\begin{equation}
  b^{\rm GMRES}_{w,s}=a_w b_{w,s}.
\end{equation}
The matrix application writes zeros to inactive output rows, so the Krylov
problem is exactly the active-row restriction of the block-diagonal system in
Equation~\eqref{eq:block-diagonal}.  In the Neumann path, inactive scratch may
be deliberately left unwritten to save bandwidth; no mathematical consumer
is allowed to read it.

\subsection{DTS VJP and child-clade propagation}

The kernel \kernel{_accumulate_gene_split_event_vjp_kernel} maps each split $n$ to
its parent activity using Equation~\eqref{eq:index-contract}.  If the parent is
inactive, it performs no atomics into child rows and no atomics into shared
parameter gradients.  Its scalar split-gradient outputs are zeroed when they
are requested.  In the staged transfer path it writes both split-side activity
bits false; when the bandwidth-saving option is enabled, this permits the
large per-species side buffers to remain uninitialized safely.

For active parents, the kernel multiplies the solved $v_i$ by normalized DTS
event probabilities and scatters those cotangents to the left and right child
rows in \code{accumulated\_rhs}.  Those atomics are precisely
Equation~\eqref{eq:child-scatter}.  Since the next earlier wave constructs its
mask only after all later-wave parents have been processed, its decision sees
the complete accumulated child cotangent.

The companion
\kernel{_accumulate_transfer_subtree_vjp_kernel} checks the same
parent-row mask before running the species-tree transfer-complement VJP and
before adding receiver-weight or child-row cotangents.  The parent activity
test and the split-side activity test are cumulative guards: either may skip
the work, and neither changes an active, nonzero contribution.

\subsection{Uninitialized scratch and persistent caches}

Some masked kernels set \code{SKIP\_INACTIVE\_SCRATCH\_ZERO=True}.  Inactive
entries of $v_i$, Neumann work arrays, or staged transfer arrays can therefore
contain arbitrary bits, including NaNs.  Correctness relies on a strict
non-read invariant:
\begin{equation}
  a_i=0 \quad\Longrightarrow\quad
  \text{every device read of row }i\text{ is masked, or the row is sanitized first}.
  \label{eq:no-read}
\end{equation}
The host uses \code{torch.where}, rather than multiplication by zero, to
sanitize $v_i$ before storing a warm start or a Hessian--vector-product cache.
This matters because $0\cdot\mathrm{NaN}=\mathrm{NaN}$.  Diagnostic residuals
are similarly replaced by zero on inactive rows.  The GMRES matrix-output path
and per-row parameter-output path explicitly zero inactive rows when their
outputs are consumed without a later mask.

\section{The separate split-side mask}
\label{sec:side-mask}

Transfer-complement VJPs stage two species vectors per gene split, one for
each child side.  Let $u_{n,q}(s)$ be the staged vector for split $n$ and side
$q\in\{L,R\}$.  The optional side mask has shape $[2N_{\rm split}]$ and is
\begin{equation}
  z_{n,q}=
  \begin{cases}
    [\max_s\abs{u_{n,q}(s)}\ne0], & \eta=0,\\
    [\sum_s\abs{u_{n,q}(s)}>\eta], & \eta>0,
  \end{cases}
  \label{eq:side-mask}
\end{equation}
where $\eta$ is \code{pibar\_side\_threshold}.  It can be produced directly
by \kernel{_accumulate_gene_split_event_vjp_kernel} or by
\kernel{_select_active_transfer_donor_sides_kernel}.  At $\eta=0$ it skips only an exactly
zero linear input, so the species-subtree correction is exactly zero.  At a
positive $\eta$ it is another threshold approximation.  This mask prunes
species-tree work for one transfer side; it does not classify a gene-clade
row and is conceptually separate from \code{USE\_ACTIVE\_MASK}.

\section{Correctness theorem and proof}

\paragraph{Proposition (exact-zero active-row pruning).}
Assume (i) clade waves are in forward topological order; (ii) rows in one wave
have no cross-row self-dependence; (iii) every $I-J_i^T$ is nonsingular; (iv)
the mask declares inactive exactly those rows with $b_i=0$; and (v) the
no-read invariant in Equation~\eqref{eq:no-read} holds.  Then the masked
reverse traversal returns exactly the same child cotangents and parameter
gradient as the unmasked traversal, up to the usual floating-point reduction
ordering.

\paragraph{Proof.}
Proceed by induction over waves in reverse topological order.  Before the root
wave, masked and unmasked traversals have the same likelihood seed.  Suppose
they have the same accumulated RHS on the current wave.  For every active row
they run the same local adjoint map and therefore produce the same $v_i$,
child scatters, and parameter contractions.  For every inactive row,
$b_i=0$; nonsingularity of $I-J_i^T$ gives the unique solution $v_i=0$, so the
unmasked child and parameter contributions are also zero.  Thus both
traversals leave identical accumulated RHS values for all earlier waves and
identical accumulated parameter gradients.  The induction reaches the leaf
waves and proves the claim. \hfill$\square$

The same argument covers exact-zero split-side pruning because its subsequent
tree transformation is linear in $u_{n,q}$.  It also explains why shared
parameters pose no problem: removing a zero summand leaves their global sum
unchanged.

\section{Positive thresholds, error, and Hessian--vector products}
\label{sec:error}

Let $P$ retain active row blocks and let $e=(I-P)b$ be the discarded RHS.  For
one row, if $\norm{J_i}\le q_i<1$ in a compatible operator norm, then
\begin{equation}
  \norm{v_i-v_i^{\rm pruned}}
  =\norm{(I-J_i^T)^{-1}e_i}
  \le \frac{\norm{e_i}}{1-q_i}.
  \label{eq:adjoint-error}
\end{equation}
If $G_i$ maps $v_i$ to the row's parameter and child cotangents, their local
error is bounded by
\begin{equation}
  \norm{G_i(v_i-v_i^{\rm pruned})}
  \le \frac{\norm{G_i}\,\norm{e_i}}{1-q_i}.
  \label{eq:gradient-error}
\end{equation}
The infinity-norm test guarantees only $\norm{e_i}_\infty<\tau$ (or
$\le\tau$, according to Equation~\eqref{eq:actual-mask}).  It does not by
itself control $1/(1-q_i)$, downstream fanout, or the total number of pruned
rows.  Threshold accuracy therefore depends on contractive, well-conditioned
local fixed points and on choosing $\tau$ relative to the desired gradient
accuracy.  The default positive threshold is a performance--accuracy policy,
not a proof of bitwise equality with the unpruned gradient.

The exact-HVP path caches the first-order $v_i$ and the same row mask, after
sanitizing inactive $v_i$ to zero.  Its tangent-adjoint solve and DTS VJP reuse
that mask.  Consequently it differentiates the algebraic adjoint with the
discrete active set held fixed.  This is the derivative of the selected
piecewise-smooth pruned computation away from a threshold crossing.  At
$r_i=\tau$ the active set can change discontinuously, and a threshold-pruned
HVP should not be interpreted as the Hessian of the smooth, fully unpruned
likelihood.  With exact-zero pruning, this distinction is normally confined
to structurally zero routes; with positive thresholds it is an explicit
approximation.

\section{Assumptions and edge cases}

\begin{itemize}
  \item A row mask is meaningful only for the wave whose local indexing it
        uses.  Passing a global-clade mask or a species mask violates
        Equation~\eqref{eq:index-contract}.
  \item $W=0$ returns an empty mask.  An all-zero root seed makes every
        downstream row exactly inactive under the strict zero criterion.
  \item NaN RHS entries are dangerous: comparisons with NaN are false, so a
        contaminated row can be classified inactive.  Finite primal and
        adjoint states are therefore a precondition, not something masking
        repairs.
  \item For $\tau=0$, the $\ge$ convention marks zero rows active, while the
        $>$ convention performs exact-zero pruning.  For $\tau>0$ the two
        conventions differ only at exact equality.
  \item The local solve is unique when $I-J_i^T$ is nonsingular.  Neumann
        convergence additionally needs a contractive self-loop; GMRES needs
        the linear system to be solvable.  Pruning does not improve a singular
        active-row system.
  \item A clade with several incoming reverse routes is tested only after
        those contributions have been summed.  It is not pruned merely
        because one parent was inactive.
  \item Floating-point atomics can change last-bit summation order.  The exact
        argument is mathematical; numerical equality is subject to ordinary
        floating-point effects and the configured finite iterative solve.
\end{itemize}

\section{Equation-to-kernel reference}

\small
\begin{longtable}{>{\raggedright\arraybackslash}p{0.39\textwidth}
                  >{\raggedright\arraybackslash}p{0.55\textwidth}}
\toprule
Kernel or host routine & Activity invariant and mathematical role \\
\midrule
\endhead
\kernel{compute_active_adjoint_row_mask} & Allocates the wave-local $[W]$ Boolean mask and selects the boundary convention in Equation~\eqref{eq:actual-mask}. \\
\kernel{_select_active_adjoint_rows_kernel} & Computes $r_w=\norm{b_w}_\infty$ from Equation~\eqref{eq:row-max}. \\
\kernel{compute_dts_forward} & Accepts the optional wave-local \code{active\_parent\_rows}; the ordinary primal leaves it absent, while backward recomputation supplies it. \\
\kernel{_reduce_single_gene_split_events_kernel} & \code{USE\_ACTIVE}: skips the one-split DTS value of an inactive parent and publishes the inactive $(-\infty,0)$ residual/offset pair. \\
\kernel{_stage_multiple_gene_split_event_reduction_kernel} & \code{USE\_ACTIVE}: skips all partial split tiles of an inactive multi-split parent. \\
\kernel{_finalize_multiple_gene_split_event_reduction_kernel} & \code{USE\_ACTIVE}: avoids reading skipped partials and leaves the inactive DTS row at $-\infty$. \\
\kernel{_prepare_reconciliation_self_loop_vjp_kernel} & \code{USE\_ACTIVE\_MASK}: forms $J_i^T$ coefficients and the initial local solve state only for active rows. \\
\kernel{_apply_reconciliation_self_loop_transpose_kernel} & \code{USE\_ACTIVE\_MASK}: applies the block of $J_i^T$ in Equation~\eqref{eq:local-adjoint}; can explicitly zero inactive outputs when a consumer requires initialized storage. \\
\kernel{_accumulate_transfer_receiver_log_probability_vjp_kernel} & \code{USE\_ACTIVE\_MASK}: accumulates the transfer-complement receiver VJP only from active $v_i$. \\
\kernel{_accumulate_reconciliation_event_vjp_kernel} & \code{USE\_ACTIVE\_MASK}: evaluates Equation~\eqref{eq:param-scatter} and writes zero per-row outputs for inactive rows. \\
\kernel{_accumulate_gene_split_event_vjp_kernel} & \code{USE\_ACTIVE\_MASK}: reads $a_{w_n}$ using \code{reduce\_idx}, implements Equation~\eqref{eq:child-scatter}, and suppresses every split and parameter scatter from an inactive parent. \\
\kernel{_select_active_transfer_donor_sides_kernel} & Builds the distinct $[2N_{\rm split}]$ mask in Equation~\eqref{eq:side-mask}. \\
\kernel{_accumulate_transfer_subtree_vjp_kernel} & Checks both parent-row and split-side activity before the linear transfer-complement subtree VJP. \\
\kernel{implicit_grad_loglik_vjp_wave} & Seeds roots, traverses waves in reverse, constructs each mask after RHS accumulation, sanitizes persistent caches, and thereby realizes the induction proof. \\
\bottomrule
\end{longtable}
\normalsize

\section{Practical interpretation}

The compile-time flags themselves do not decide which biology is relevant;
they merely remove a runtime branch and its guarded arithmetic from a Triton
specialization when no mask was supplied.  The mathematical decision is the
wave-local projection in Equation~\eqref{eq:actual-mask}.  Exact-zero rows can
be removed because the adjoint solve and every outgoing VJP are linear in the
incoming cotangent.  Positive-threshold rows can be removed because the user
has elected to approximate those small cotangents as zero.  Reverse wave order,
consistent guarding of every consumer, and sanitation of any persistent
scratch are what turn that row decision into safe subclade pruning.

\end{document}
